\section{The geometric side of the metaplectic trace formula}

We follow [A2,A3].

\subsection{The kernel function}

Let ${\cal{C}}$ denote the set of equivalence classes of
$\gamma \in G({\mathbb{Q}})$ as in [A2, section 2].
Recall that ${\cal{C}}$ is in bijection
with the set of equivalence classes of pairs $(M,{\cal{O}}_M)$,
where $P=MN $ is a standard parabolic and
${\cal{O}}_M\subset M({\mathbb{Q}})$ is a
conjugacy class of elliptic elements.

For ${\cal{O}}\in {\cal{C}}$, let 
\[
K_{{\cal{O}}}(x,y)=\sum_{\gamma \in {\cal{O}}}f(x\gamma y^{-1}),
\]
and more generally, let 
\begin{equation}
K_{P,{\cal{O}}}(x,y)=\sum_{\gamma \in M({\mathbb{Q}})\cap {\cal{O}}}
\int_{N({\mathbb{A}})}f(x\gamma y^{-1})dn,
\end{equation}
where $P$ is a parabolic as above.

\begin{lemma}
For $\overline{K_0}$-finite
$f\in {\cal{H}}(\overline{G({\mathbb{A}})})$, the sum
$K_P(x,y)=\sum_{{\cal{O}}\in {\cal{C}}}K_{{\cal{O}}}(x,y)$
converges to the kernel
function of $R_P$, where $R_P$ denotes the right
regular representation of
${\cal{H}}(\overline{G({\mathbb{A}})})_{gen}$ on
$L^2(N({\mathbb{A}})M({\mathbb{Q}})\backslash \overline{G({\mathbb{A}})})_{gen}$.
\end{lemma}

proof: For
$\phi \in C_c^\infty (N({\mathbb{A}})M({\mathbb{Q}})
\backslash \overline{G({\mathbb{A}})})_{gen}$, we have 
\begin{eqnarray*}
(R_P(f)\phi )(x) &=&\int_{\overline{G({\mathbb{A}})}}\phi (g)f(xg^{-1})dg \\
&=&\int_{\overline{G({\mathbb{A}})}}\phi (n^{-1}\gamma ^{-1}g)f(xg^{-1})dg \\
&=&\int_{\overline{G({\mathbb{A}})}}\phi (g)f(xn\gamma g^{-1})dg \\
&=&\int_{M({\mathbb{Q}})\backslash \overline{G({\mathbb{A}})}}\phi (g)\sum_{\gamma \in
M({\mathbb{Q}})}f(xn\gamma g^{-1})dg \\
&=&\int_{N({\mathbb{A}})M({\mathbb{Q}})\backslash \overline{G({\mathbb{A}})}}\phi
(g)\sum_{\gamma \in M({\mathbb{Q}})}\int_{N({\mathbb{A}})}f(xn\gamma g^{-1})dndg.
\end{eqnarray*}

The convergence of the inner sum is established, for example, by passing to
the place at infinity and carrying out the estimates in the usual manner.
For the description of
$G({\mathbb{Q}})\cap (\overline{G({\mathbb{R}})\times K^f})$,
see [K2] for example.$\Box $

\subsubsection{Some analytic reduction theory}

Let 
\[
N({\mathbb{A}})_u=\{\left( 
\begin{array}{cc}
1 & n \\ 
0 & 1
\end{array}
\right) \;|\;|n|_{{\mathbb{A}}}\leq u\}. 
\]
For $\omega \subset N({\mathbb{A}})_u\cdot A({\mathbb{A}})^1$
relatively compact, $T_0\in {\mathfrak{a}}^{+}$, let 
\[
S(T_0,\omega )=
\{nak\in \overline{\omega \cdot A({\mathbb{R}})^0\cdot K_0}
\;|\;\alpha (H(a)-T_0)>0\}. 
\]
This is a Siegel domain of $\overline{G({\mathbb{A}})}$. It has finite volume
with respect to the measure $\delta (a)dndadk$. For $T_0\in {\mathfrak{a}}^{+}$
and $\omega \subset N({\mathbb{A}})_u\cdot A({\mathbb{A}})^1$ relatively
compact both
``sufficiently large'' ($u\geq 1/2,\;T_0>\sqrt{3}/2$), we have 
\[
\overline{G({\mathbb{A}})}=G({\mathbb{Q}})\cdot S(T_0,\omega ). 
\]
In particular, $G({\mathbb{Q}})\backslash \overline{G({\mathbb{A}})}$ is
finite volume.

\begin{lemma}
There exists an $N>0$ such that for all $g\in \overline{G({\mathbb{A}})}$,
$T_0\in {\mathfrak{a}}^{+}$, we have 
\[
\sum_{\delta \in B({\mathbb{Q}})\backslash G({\mathbb{Q}})}
\tau _B^{\wedge }(H(\delta x)-T_0)\leq 1.
\]
\end{lemma}

This is an immediate consequence of the corresponding result in [A2].

For $T_0>T_1$ in ${\mathfrak{a}}^{+}$, let 
\[
S(T_0,T_1,\omega )=\{nak\in S(T_0,\omega )\;|\;\omega (H(a)-T_0)\leq 0\}. 
\]
Let $F^G(x,T_1,T_0)$ denote the characteristic function of the image
$\overline{S(T_0,T_1,\omega )}\subset
G({\mathbb{Q}})\backslash \overline{G({\mathbb{A}})}$ of
$S(T_0,T_1,\omega )$ under the quotient map
$\overline{G({\mathbb{A}})}\rightarrow
G({\mathbb{Q}})\backslash \overline{G({\mathbb{A}})}$. Let
$F^B(x,T_1,T_0)$
denote the characteristic function of the set of all
$xB({\mathbb{Q}})\in B({\mathbb{Q}})\backslash \overline{G({\mathbb{A}})}$ such
that $\delta x\in \overline{\omega \cdot A({\mathbb{R}})^o\cdot K_0}$
for some $\delta \in B({\mathbb{Q}})$.

\begin{lemma}
For $\omega \subset N({\mathbb{A}})_u\cdot A({\mathbb{A}})^1$ relatively
compact and
$T_0\in {\mathfrak{a}}^{+}$ sufficiently large,
$x\in \overline{G({\mathbb{A}})}$, we
have 
\[
\sum_{P\in \{B,G\}}
\sum_{\delta \in P({\mathbb{Q}})\backslash G({\mathbb{Q}})}
F^P(\delta x,T_1,T_0)\tau _P(H(\delta x)-T_1)=1.
\]
\end{lemma}

This is an immediate consequence of the corresponding result in [A2].

\subsubsection{The truncated kernel}

The auxiliary kernel function is by definition 
\[
k^T(x,f)=\sum_{P\in \{B,G\}}(-1)^{\dim (A_P)}
\sum_{\delta \in P({\mathbb{Q}})\backslash G({\mathbb{Q}})}
K_P(\delta x,\delta x)\tau _P^{\wedge }(H(\delta x)-T). 
\]
This sum is finite by Lemma 4.2 above. Using the previous Lemma 4.3,
this may be rewritten as 
\begin{eqnarray*}
k^T(x,f) &=&\sum_{P,Q\in \{B,G\}}(-1)^{\dim (A_P)}\times \\
&&\;\;\;\;\;\;\;\;\;\;\;\times
\sum_{\delta \in Q({\mathbb{Q}})\backslash G({\mathbb{Q}})}
K_Q(\delta x,\delta x)\tau _Q^{\wedge }(H(\delta x)-T)
F^P(\delta x,T,T_0)\tau _P^Q(H(\delta x)-T),
\end{eqnarray*}
for $T>T_0$ regular. This may be rewritten as
\begin{eqnarray}
k^T(x,f) &=&K_G(x,x)F^G(x,T,T_0)+ \\
&&\sum_{\delta \in B({\mathbb{Q}})\backslash G({\mathbb{Q}})}
\tau _B(H(\delta x)-T)
F^B(\delta x,T,T_0)(K_G(\delta x,\delta x)-K_B(\delta x,\delta x)). 
\nonumber
\end{eqnarray}

We want to prove

\begin{lemma}
$\int_{G({\mathbb{Q}})\backslash \overline{G({\mathbb{A}})}}|k^T(x,f)|dx<\infty .$
\end{lemma}

We majorize this by the sum of the integrals of the absolute values of the
two terms on the right hand side of (4.2). The integral over the first term
on the right hand side of (4.2), 
\[
\int_{G({\mathbb{Q}})\backslash \overline{G({\mathbb{A}})}}|K_G(x,x)F^G(x,T,T_0)|dx 
\]
clearly converges since $F^G(x,T,T_0)$ is compactly supported. To estimate
the second term, we first write that 
\begin{eqnarray}
&&\int_{G({\mathbb{Q}})\backslash \overline{G({\mathbb{A}})}}|
\sum_{\delta \in B({\mathbb{Q}})\backslash G({\mathbb{Q}})}
\tau _B(H(\delta x)-T)\times \\
&&\;\;\;\;\;\;\;\;\;\;\;\;\;\;\;\;\;\;\;\;\times F^B(\delta
x,T,T_0)(K_G(\delta x,\delta x)-K_B(\delta x,\delta x))|dx \\
&<&<\int_{G({\mathbb{Q}})\backslash \overline{G({\mathbb{A}})}}
\tau _B(H(\delta x)-T)F^B(\delta x,T,T_0)\times  \nonumber \\
&&\;\;\;\;\;\;\;\;\;\;\;\;\;\;\;\;\;\;\;\;\;\;\;\times
|\sum_{\delta \in G({\mathbb{Q}})}f(x^{-1}\delta x)
-\int_{N({\mathbb{A}})}f(x^{-1}\mu n^{\prime }x)dn^{\prime }|dx.
\end{eqnarray}
Next, we again pass to the archimedian place. Let 
\begin{eqnarray}
\Gamma _0 &=&N({\mathbb{Q}})\cap (\overline{G({\mathbb{R}})\times K^f}),
\;\Gamma_\infty
=N({\mathbb{R}})\cap (\overline{G({\mathbb{R}})\times K_0^f}),\; \\
A({\mathbb{R}})_T &=&\{\prod_{v|\infty}
\left( 
\begin{array}{cc}
t_v & 0 \\ 
0 & t_v^{-1}
\end{array}
\right) \in A({\mathbb{R}})\;|\;|t_v|>T^{1/2}\},
\;G({\mathbb{R}})_T=N({\mathbb{R}})\cdot A({\mathbb{R}})_T\cdot K_\infty ,  
\nonumber
\end{eqnarray}
so $\Gamma _0\subset \Gamma _\infty $ is a co-compact discrete abelian
group. The above upper bound (1.10) is finite if 
\begin{equation}
\int_{\Gamma _0\backslash \overline{G({\mathbb{R}})_T}}|\sum_{\gamma \in \Gamma
_0}f(x^{-1}\gamma x)-\int_{\Gamma _\infty }f(x^{-1}nx)dn|dx<\infty .
\end{equation}
We may apply the theory of Fourier transforms, as in [W1], to the 
quotient $\Gamma _0\backslash \Gamma _\infty $. 
Choosing Haar measures appropriately,
there is a non-trivial character 
$\psi $ of $\Gamma _\infty $ such that the Fourier transform 
\[
\varphi ^{\wedge }(\gamma )
=\int_{\Gamma _\infty }\varphi (\gamma ^{\prime})
\psi (<\gamma ,\gamma ^{\prime }>)d\gamma ^{\prime },
\;\;\;\varphi \in L^2(\Gamma _0\backslash \Gamma _\infty ), 
\]
satisfies 
\[
\varphi ^{\wedge }(1)=\int_{\Gamma _\infty }\varphi (\gamma ^{\prime
})d\gamma ^{\prime }. 
\]

Take $\varphi (\gamma )=\varphi _x(\gamma )=f(x^{-1}\gamma x)$, so by the
Poisson summation formula [W1], 
\[
\sum_{\gamma \in \Gamma _0}f(x^{-1}\gamma x)-\int_{\Gamma _\infty
}f(x^{-1}nx)dn=\sum_{\gamma ^{*}\in \Gamma _0^{*}}\varphi _x^{\wedge
}(\gamma ^{*}), 
\]
where $\Gamma _0^{*}$ denotes the dual lattice of $\Gamma _0$. Since $f$ is
smooth, $\varphi _x^{\wedge }$ is rapidly decreasing. For sufficiently
regular $T$, standard arguments now suffice to show that, for any constant $%
C>0$, we have 
\[
\int_{\Gamma _0\backslash \overline{G({\mathbb{R}})_T}}|\sum_{\gamma ^{*}\in
\Gamma _0^{*}}\varphi _x^{\wedge }(\gamma ^{*})|dx<<_{C,f,n}T^{-C}. 
\]

This verifies (1.12), which proves the lemma. $\Box $

The above proof actually establishes more, namely,

\begin{corollary}
(of the above proof) For each constant $C>0$, we have 
\[
\int_{G({\mathbb{Q}})\backslash \overline{G({\mathbb{A}})}}k^T(x,f)dx-\int_{G({\mathbb{Q}}%
)\backslash \overline{G({\mathbb{A}})_T}}K^G(x,x)dx<<_{C,f,n}T^{-C}. 
\]
Here 
\[
G({\mathbb{A}})_T=N({\mathbb{A}})_u\cdot A({\mathbb{A}})^T\cdot K 
\]
where $A({\mathbb{A}})^T=\{a\in A({\mathbb{A}})\;|\;|\delta (a)|_{\mathbb A}\leq T\}$.
\end{corollary}

In fact, the proof yields more than this. If we let 
\[
k_{{\cal{O}}}^T(x,f)=\sum_{P\in \{B,G\}}(-1)^{\dim (A_P)}\sum_{\delta \in
P({\mathbb{Q}})\backslash G({\mathbb{Q}})}K_{P,{\cal{O}}}(\delta x,\delta x)\tau
_P^{\wedge }(H(\delta x)-T), 
\]
for ${\cal{O}}\in {\cal{C}}$, then the above proof establishes

\begin{corollary}
(of the above proof) For each constant $C>0$, we have 
\[
\sum_{{\cal{O}}\in {\cal{C}}}\int_{G({\mathbb{Q}})\backslash \overline{G(\mathbb{%
A})}}|k_{{\cal{O}}}^T(x,f)dx-F^G(x,T,T_0)K_{G,{\cal{O}}%
}(x,x)|dx<<_{C,f,n}T^{-C}. 
\]
\end{corollary}

\subsection{Computing the integral of $k^T(x,f)$}

For a class ${\cal{O}}\in {\cal{C}}$ and a standard parabolic $P=MN$,
let 
\[
J_{P,{\cal{O}}}(x,y)=\sum_{\mu \in M({\mathbb{Q}})\cap {\cal{O}}%
}\sum_{\varsigma \in N_{\mu _s}({\mathbb{Q}})\backslash N({\mathbb{Q}})}\int_{N_{\mu
_s}({\mathbb{A}})}f(x^{-1}\varsigma ^{-1}\mu n\varsigma y)dn, 
\]
where $\mu =\mu _s\mu _u$ denotes the Jordan decomposition (over $\overline{%
{\mathbb{Q}}}$), and let 
\[
j_{{\cal{O}}}^T(x,f)=\sum_{P\in \{B,G\}}(-1)^{\dim (A_P)}\sum_{\delta \in
P({\mathbb{Q}})\backslash G({\mathbb{Q}})}J_{P,{\cal{O}}}(\delta x,\delta x)\tau
_P^{\wedge }(H(\delta x)-T). 
\]
the above proof establishes

\begin{corollary}
(of the above proof) For each constant $C>0$, we have 
\[
\sum_{{\cal{O}}\in {\cal{C}}}\int_{G({\mathbb{Q}})\backslash \overline{G(\mathbb{%
A})}}|j_{{\cal{O}}}^T(x,f)dx-F^G(x,T,T_0)J_{G,{\cal{O}}%
}(x,x)|dx<<_{C,f,n}T^{-C}. 
\]
\end{corollary}

Let 
\begin{eqnarray*}
J_{{\cal{O}}}^T(f) &=&\int_{G({\mathbb{Q}})\backslash \overline{G({\mathbb{A}})}}j_{%
{\cal{O}}}^T(x,f)dx, \\
K_{{\cal{O}}}^T(f) &=&\int_{G({\mathbb{Q}})\backslash \overline{G({\mathbb{A}})}}k_{%
{\cal{O}}}^T(x,f)dx.
\end{eqnarray*}

\begin{proposition}
Both $
J_{{\cal{O}}}^T(f)
$ and $K_{{\cal{O}}}^T(f)$ are polynomials in
$T$.
\end{proposition}

Define
\begin{equation}
J_{{\cal{O}}}(f)
=J_{{\cal{O}}}^T(f)|_{T=0}.
\end{equation}

We will prove this below. First, we record a

\begin{corollary}
$\sum_{{\cal{O}}\in {\cal{C}}}J_{{\cal{O}}}^T(f)=\sum_{{\cal{O}}\in 
{\cal{C}}}K_{{\cal{O}}}^T(f)$.
\end{corollary}

proof: The difference between them is a polynomial, $p$ say, in $T$ by the
proposition. Since $J_{G,{\cal{O}}}(x,y)=K_{G,{\cal{O}}}(x,y)$, this
polynomial $p$ must be the zero polynomial by corollaries 4.6 
and 4.7. $\Box $

Let ${\cal{O}}\in {\cal{C}}$ denote an unramified semi-simple class [A2,
p. 949]. If $\mu \in {\cal{O}}$ then its centralizer $G_\mu $ is either $G$
or a Cartan subgroup of $G$ defined over ${\mathbb{Q}}$.

As in [A2], we have 
\begin{equation}
j_{{\cal{O}}}^T(x,f)=
\sum_{\eta \in G_\mu ({\mathbb{Q}})\backslash G({\mathbb{Q}})}
f(x^{-1}\eta ^{-1}\mu \eta x)\sum_{P\in \{B,G\}}(-1)^{\dim (A_P)}
\tau_P^{\wedge }(H(\eta x)-T),
\end{equation}
so 
\begin{eqnarray}
J_{{\cal{O}}}^T(f)
&=&\int_{G_\mu ({\mathbb{Q}})\backslash \overline{G({\mathbb{A}})}}
f(x^{-1}\mu x)\sum_{P\in \{B,G\}}(-1)^{\dim (A_P)}
\tau _P^{\wedge}(H(x)-T)dx \\
&=&vol(A_{{\cal{O}}}({\mathbb{R}})^0G_\mu ({\mathbb{Q}})
\backslash \overline{G_\mu ({\mathbb{A}})}
\int_{\overline{G_\mu ({\mathbb{A}})}\backslash \overline{G({\mathbb{A}})}}
f(x^{-1}\mu x)v_{{\cal{O}}}(x,T)dx,  \nonumber
\end{eqnarray}
where $A_{{\cal{O}}}$ denotes the center of $M_\mu $ and 
\[
v_{{\cal{O}}}(x,T)
=\int_{A_{{\cal{O}}}({\mathbb{R}})^0}\sum_{P\in
\{B,G\}}(-1)^{\dim (A_P)}\tau _P^{\wedge }(H(ax)-T)da. 
\]
Note that the global cocycle $\beta $ splits over
$A_{{\cal{O}}}({\mathbb{R}})^o\subset A({\mathbb{R}})^o$.
In the case where $A_{\cal{O}}=A$, this may be
computed to be 
\[
v_{{\cal{O}}}(x,T)=2(T-H(x)). 
\]

If $\mu \in {\cal{O}}$ is semi-simple but not regular then 
\[
j_{{\cal{O}}}^T(x,f)
=f(\mu )\sum_{P\in \{B,G\}}(-1)^{\dim (A_P)}
\tau_P^{\wedge }(H(x)-T) 
\]
so 
\begin{equation}
J_{{\cal{O}}}^T(f)=f(\mu )
\int_{G({\mathbb{Q}})\backslash \overline{G({\mathbb{A}})}}
\sum_{P\in \{B,G\}}(-1)^{\dim (A_P)}
\tau _P^{\wedge }(H(x)-T)dx.  \nonumber
\end{equation}


